% Part: normal-modal-logic
% Chapter: filtrations
% Section: examples-of-filtrations

\documentclass[../../../include/open-logic-section]{subfiles}

\begin{document}

\olfileid[pt]{nml}{fil}{exf}

\olsection{Exemplos de filtrações}

Ainda não mostramos que existe alguma filtração. Na verdade, para todo
modelo $\mModel{M}$ há muitas filtrações de $\mModel{M}$ por meio de
$\Gamma$. Destacamos duas: a filtração mais fina e a mais grossa.
Filtrações dos mesmos modelos diferem em sua relação de acessibilidade
(pois \olref[fil]{defn:filtration} estipula diretamente como devem ser
$W^*$ e~$V^*$). A filtração mais fina terá o menor número possível de
mundos relacionados, ao passo que a mais grossa terá o maior número
possível.

\begin{defn}
  Se $\Gamma$ é fechado sob subfórmulas, a filtração \emph{mais fina}
  $\mModel{M^*}$ de um modelo $\mModel{M}$ é definida estipulando:
  \[
  R^*[u][v] \quad \text{se, e somente se,} \quad \exists u'\in [u] \;
  \exists v' \in [v] : Ru'v'.
  \]
\end{defn}

\begin{prop}\ollabel{prop:finest}
  A filtração mais fina $\mModel{M^*}$ é de fato uma filtração. 
\end{prop}

\begin{proof}
  Precisamos verificar que $R^*$, assim definida, satisfaz
  \olref[fil]{defn:filtration}\olref[fil]{defn:filtration-R}.
  Verificamos sucessivamente as três condições.

  Se $Ruv$, então, como $u \in [u]$ e $v \in [v]$, também $R^*[u][v]$;
  portanto, \olref[fil]{defn:filtration-R1} é satisfeita.
  
  \iftag{prvBox}{\iftag{probBox}{Deixamos a verificação de
      \olref[fil]{defn:filtration-R2} como exercício.}{Para
        \olref[fil]{defn:filtration-R2}, suponha $\Box!A \in \Gamma$,
        $R^*[u][v]$ e $\mSat{M}{\Box!A}[u]$. Pela definição de
        $R^*$, existem $u' \equiv u$ e $v' \equiv v$ tais que
        $Ru'v'$. Como $u$ e $u'$ concordam em $\Gamma$, também
        $\mSat{M}{\Box!A}[u']$; logo, $\mSat{M}{!A}[v']$. Pelo fechamento
        de $\Gamma$ sob sub-!!{formula}s, $v$ e $v'$ concordam em
        $!A$; portanto, $\mSat{M}{!A}[v]$, como desejado.}}{}
  
  \iftag{prvDiamond}{\iftag{probDiamond}{Deixamos a verificação de
      \olref[fil]{defn:filtration-R3} como exercício.}{Para verificar
        \olref[fil]{defn:filtration-R3}, suponha $\Diamond!A \in
        \Gamma$, $R^*[u][v]$ e $\mSat{M}{!A}[v]$. Pela definição de
        $R^*$, existem $u' \equiv u$ e $v' \equiv v$ tais que
        $Ru'v'$. Como $v$ e $v'$ concordam em~$\Gamma$ e $\Gamma$ é
        fechado sob sub-!!{formula}s, também $\mSat{M}{!A}[v']$; logo,
        $\mSat{M}{\Diamond !A}[u']$. Como $u$ e $u'$ também concordam
        em $\Gamma$, $\mSat{M}{\Diamond !A}[u]$.}}{}
\end{proof}

\begin{probtag}{probBox,probDiamond}
  Complete a demonstração de \olref[nml][fil][exf]{prop:finest}.
\end{probtag}

\begin{defn}
  Se $\Gamma$ é fechado sob subfórmulas, a filtração \emph{mais grossa}
  ~$\mModel{M^*}$ de um modelo~$\mModel{M}$ é definida estipulando que
  $R^*[u][v]$ se, e somente se,
  \iftag{notprvBox,notprvDiamond}{a condição a seguir é
    satisfeita:}{\emph{ambas} as condições a seguir são satisfeitas:}
  \begin{tagenumerate}{prvBox,prvDiamond}
  \tagitem{prvBox}{\ollabel{defn:coarsest-Box}Se $\Box!A \in \Gamma$
    e $\mSat{M}{\Box!A}[u]$, então
    $\mSat{M}{!A}[v]$\iftag{prvDiamond}{;}{.}}{}
  \tagitem{prvDiamond}{\ollabel{defn:coarsest-Diamond}Se $\Diamond!A
    \in \Gamma$ e $\mSat{M}{!A}[v]$, então
    $\mSat{M}{\Diamond!A}[u]$.}{}
  \end{tagenumerate}
\end{defn}

\begin{prop}
  A filtração mais grossa $\mModel{M^*}$ é de fato uma filtração. 
\end{prop}

\begin{proof}
  Dada a definição de $R^*$, a única condição que resta verificar é a
  implicação de $Ruv$ para $R^*[u][v]$. Suponha, portanto, $Ruv$.
  \iftag{prvBox}{Suponha $\Box!A \in \Gamma$ e
    $\mSat{M}{\Box!A}[u]$; então, evidentemente,
    $\mSat{M}{!A}[v]$\iftag{prvDiamond}{, e
      \olref{defn:coarsest-Box} é
      satisfeita. }{.}}{}\iftag{prvDiamond}{Suponha $\Diamond!A \in
    \Gamma$ e $\mSat{M}{!A}[v]$. Então $\mSat{M}{\Diamond!A}[u]$,
    pois~$Ruv$\iftag{prvBox}{, e \olref{defn:coarsest-Diamond} é
      satisfeita}{}.}{}
\end{proof}

\begin{ex}
  Seja $W = \PosInt$, valha $Rnm$ sse $m = n + 1$ e seja
  $V(p) = \Setabs{2n}{n \in \Nat}$. O modelo $\mModel{M} =
  \tuple{W, R, V}$ é representado em \olref{fig:ex-filtration}. Os
  mundos são $1$, $2$ etc.; cada mundo acessa exatamente um outro
  mundo---seu sucessor---e $p$ é verdadeira em todos, e somente, os
  números pares.

  \begin{figure}
    \centering
    \begin{tikzpicture}[modal]
      \node[world] (1) [label=below:\mFalse{p}] {$1$};
      \node[world] (2) [label=below:\mTrue{p}, right=of 1] {$2$};
      \node[world] (3) [label=below:\mFalse{p}, right=of 2] {$3$};
      \node[world] (4) [label=below:\mTrue{p}, right=of 3] {$4$};
      \node[phantom] (5) [right=of 4] {};
      \draw[->] (1) to (2);
      \draw[->] (2) to (3);
      \draw[->] (3) to (4);
      \draw[dotted] (4) to (5);
    \end{tikzpicture}

    \begin{tikzpicture}[modal]
      \node[world] (1) [label=below:\mFalse{p}] {$[1]$};
      \node[world] (2) [label=below:\mTrue{p}, right=of 1] {$[2]$};
      \draw[->,bend left] (1) to (2);
      \draw[->,bend left] (2) to (1);
      
      \node[world] (3) [label=below:\mFalse{p},right=of 2] {$[1]$};
      \node[world] (4) [label=below:\mTrue{p}, right=of 3] {$[2]$};
      \draw[->,bend left] (3) to (4);
      \draw[->,bend left] (4) to (3);
      \draw[->,reflexive above] (4) to (4);
    \end{tikzpicture}
    \caption{Um modelo infinito e suas filtrações.}
    \ollabel{fig:ex-filtration}
  \end{figure}
  
  Seja agora $\Gamma$ o conjunto das sub-!!{formula}s de~$\Box p \lif p$,
  O.s.\@, $\{p, \Box p, \Box p \lif p\}$. $p$ é verdadeira em todos,
  e somente, os números pares; $\Box p$ é verdadeira em todos, e somente,
  os números ímpares; logo, $\Box p \lif p$ é verdadeira em todos, e
  somente, os números pares. Em outras palavras, todo número ímpar torna
  $\Box p$ verdadeira, mas $p$ e~$\Box p \lif p$ falsas; todo número par
  torna $p$ e $\Box p \lif p$ verdadeiras mas $\Box p$ falsa. Portanto,
  $W^* = \{ [1], [2] \}$, em que $[1] = \{1, 3, 5, \dots\}$ e
  $[2] = \{2, 4, 6, \dots\}$. Como $2 \in V(p)$, $[2] \in V^*(p)$;
  como $1 \notin V(p)$, $[1] \notin V^*(p)$. Logo, $V^*(p) = \{[2]\}$.

  Toda filtração baseada em $W^*$ deve ter uma relação de acessibilidade
  que inclua $\tuple{[1], [2]}, \tuple{[2],[1]}$: como $R12$, devemos ter
  $R^*[1][2]$ por \olref[fil]{defn:filtration}
  \olref[fil]{defn:filtration-R1}; como $R23$, devemos ter $R^*[2][3]$,
  e $[3]=[1]$. Ela não pode incluir $\tuple{[1],[1]}$: se incluísse,
  teríamos $R^*[1][1]$, $\mSat{M}{\Box p}[1]$, mas $\mSat/{M}{p}[1]$,
  em contradição com \olref[fil]{defn:filtration-R2}. Nada exige nem
  exclui $R^*[2][2]$. Há, portanto, duas filtrações possíveis de
  ~$\mModel{M}$, correspondentes às duas relações de acessibilidade
  \[
  \{\tuple{[1],[2]}, \tuple{[2],[1]}\} \text{ e }
  \{\tuple{[1],[2]}, \tuple{[2],[1]}, \tuple{[2],[2]}\}.
  \]
  Em ambos os casos, $p$ e $\Box p \lif p$ são falsas e $\Box p$ é
  verdadeira em~$[1]$; $p$ e $\Box p \lif p$ são verdadeiras e
  $\Box p$ é falsa em~$[2]$.
\end{ex}


\begin{prob}
  Considere o modelo $\mModel{M} = \tuple{W, R, V}$ a seguir, em que
  $W = \Setabs{0\sigma}{\sigma \in \Bin^*}$ é o conjunto das sequências
  de $0$s e~$1$s que começam por~$0$, com $R\sigma\sigma'$ sse
  $\sigma' = \sigma 0$ ou $\sigma' = \sigma 1$, e
  $V(p) = \Setabs{\sigma 0}{\sigma \in \Bin^*}$ e
  $V(q) = \Setabs{\sigma 1}{\sigma \in \Bin^* \setminus \{1\}}$.
  Eis uma representação:
  \begin{center}
    \begin{tikzpicture}[modal,
        node distance=2cm
      ]

      \node[world] (0) [label={below:\mTrue{p}\\\mFalse{q}},font=\tiny] {$0$};
      \node[world] (00) [label={below:\mTrue{p}\\\mFalse{q}},
        above right=of 0,font=\tiny] {$00$};
      \node[world] (000) [label={below:\mTrue{p}\\\mFalse{q}},
        right=of 00, yshift=.8cm, font=\tiny] {$000$};
      \node[phantom] (0000) [right=of 000, yshift=.5cm] {};
      \node[phantom] (0001) [right=of 000, yshift=-.5cm] {};
      \node[world] (001) [label={below:\mFalse{p}\\\mTrue{q}},
        right=of 00, yshift=-.8cm, font=\tiny] {$001$};
      \node[phantom](0010) [right=of 001, yshift=.5cm] {};
      \node[phantom](0011) [right=of 001, yshift=-.5cm] {};
      \node[world] (01) [label={below:\mFalse{p}\\\mTrue{q}},
        below right=of 0,font=\tiny] {$01$};
      \node[world] (010) [label={below:\mTrue{p}\\\mFalse{q}},
        right=of 01, yshift=.8cm,font=\tiny] {$010$};
      \node[phantom] (0100) [right=of 010, yshift=.5cm] {};
      \node[phantom] (0101) [right=of 010, yshift=-.5cm] {};
      \node[world] (011) [label={below:\mFalse{p}\\\mTrue{q}},
        right=of 01, yshift=-.8cm,font=\tiny] {$011$};
      \node[phantom] (0110) [right=of 011, yshift=.5cm] {};
      \node[phantom] (0111) [right=of 011, yshift=-.5cm] {};

      \draw[->] (0) to (00);
      \draw[->] (0) to (01);
      \draw[->] (00) to (000);
      \draw[dotted] (000) to (0000);
      \draw[dotted] (000) to (0001);
      \draw[->] (00) to (001);
      \draw[dotted] (001) to (0010);
      \draw[dotted] (001) to (0011);
      \draw[->] (01) to (010);
      \draw[dotted] (010) to (0100);
      \draw[dotted] (010) to (0101);
      \draw[->] (01) to (011);
      \draw[dotted] (011) to (0110);
      \draw[dotted] (011) to (0111);
    \end{tikzpicture}
  \end{center}
  Temos $\mSat/{M}{\Box(p \lor q) \lif (\Box p \lor \Box q)}[w]$ para
  todo~$w$.

  Seja $\Gamma$ o conjunto das sub-!!{formula}s de~$\Box(p \lor q) \lif
  (\Box p \lor \Box q)$. Quais são $W^*$ e~$V^*$? Qual é a relação de
  acessibilidade da filtração mais fina de~$\mModel{M}$? E da mais grossa?
\end{prob}

\end{document}
